%% IEEE TAES manuscript draft
%% Build: pdflatex twice, then bibtex is not used (bibitems inline).
%% Requires IEEEtran.cls (https://www.michaelshell.org/tex/ieeetran/).
\documentclass[journal]{IEEEtran}

\usepackage{amsmath,amssymb,amsfonts}
\usepackage{graphicx}
\usepackage{booktabs}
\usepackage{array}
\usepackage[hidelinks]{hyperref}
\usepackage{siunitx}

\newcommand{\nT}{\,\mathrm{nT}}
\newcommand{\bx}{\mathbf{x}}
\newcommand{\bA}{\mathbf{A}}
\newcommand{\bH}{\mathbf{H}}
\newcommand{\bF}{\boldsymbol{\Phi}}

\begin{document}

\title{Factor-Graph Aeromagnetic Compensation for Airborne\\
Magnetic-Anomaly Navigation: A Neural-Network-Free\\
Sliding-Window Alternative to Online EKF Calibration}

\author{\IEEEauthorblockN{Anonymous}%
\thanks{Manuscript draft. Code and reproducible experiments:
MagNav.jl research branch \texttt{claude/fgo-problem-research-bllgbr}.}}

\markboth{IEEE Transactions on Aerospace and Electronic Systems (draft)}%
{Factor-Graph Aeromagnetic Compensation for Magnetic-Anomaly Navigation}

\maketitle

\begin{abstract}
Airborne magnetic-anomaly navigation (MagNav) matches a scalar magnetometer
against a magnetic-anomaly map to bound inertial drift in GNSS-denied flight. Its
central obstacle is aeromagnetic \emph{interference}: the aircraft's own field
corrupts the scalar measurement and must be compensated online, because a
dedicated calibration flight is not always available (the ``cold-start'' regime).
Recent work augments an extended Kalman filter (EKF) with online Tolles--Lawson
(TL) coefficients \emph{and} a residual neural network (NN), stabilized by a
specialized training scheme. We show that a factor-graph optimization (FGO)
formulation attains the same accuracy \emph{without any neural network}. We carry
the TL compensation coefficients as factor-graph variables and estimate them
jointly with navigation in a fixed-lag sliding-window smoother, so the calibration
adapts to time-varying interference along the flight. On the primary evaluation
line of the SGL~2020 dataset the sliding-window FGO reaches
\SI{32.6}{\meter}/\SI{14.2}{\meter} distance-RMS (DRMS) on the two uncompensated
cabin magnetometers, matching or beating a reproduced online EKF+TL+NN
(\SI{40.0}{}/\SI{17.5}{\meter}) and the values published for it
(\SI{37}{}/\SI{14}{\meter}). Across five long survey lines from three flights and
two maps, the sliding-window FGO beats a causal online-TL EKF on
\num{8} of \num{9} cold-start cases; crucially, the causal EKF \emph{diverges} to
tens of kilometers on the noisy magnetometer while the batch/window smoother stays
bounded, because its per-window relinearization lets calibration that only becomes
observable later protect earlier navigation. We also cast physical
scalar-magnetometer error terms as graph variables and analyze the observability
of joint compensation and navigation.
\end{abstract}

\begin{IEEEkeywords}
Magnetic-anomaly navigation, aeromagnetic compensation, Tolles--Lawson, factor
graph optimization, fixed-lag smoothing, GNSS-denied navigation.
\end{IEEEkeywords}

\section{Introduction}
\IEEEPARstart{M}{agnetic}-anomaly navigation estimates aircraft position by
matching a measured scalar magnetic field against a surveyed crustal-anomaly map,
providing an absolute, unjammable position reference that bounds the drift of an
inertial navigation system (INS) when GNSS is denied~\cite{canciani2017}. The
enabling public benchmark is the Signal Enhancement for Magnetic Navigation (SGL
2020) challenge dataset~\cite{gnadt2020}, on which most recent work is evaluated.

The dominant error source is not the map but the aircraft: permanent and induced
magnetization and eddy-current fields perturb the scalar measurement by hundreds
of nanotesla, dwarfing the anomaly signal ($\sim$tens of \nT). Classical
aeromagnetic compensation uses the Tolles--Lawson (TL) model~\cite{tolles1950}, a
linear regression of the interference onto attitude-dependent terms built from a
three-axis fluxgate, whose coefficients are fit on a dedicated high-altitude
calibration flight. When no calibration flight is available---the operationally
important \emph{cold-start} regime---the coefficients must be learned online,
jointly with navigation.

Hager \emph{et al.}~\cite{hager2026} address the cold start by augmenting an EKF
with online TL coefficients and a residual neural network (NN), the latter
learned online and stabilized by a dedicated training design; this is the state
of the art we compare against. In this paper we ask whether the neural network is
necessary. We show it is not: a factor-graph formulation that treats the TL
coefficients as time-varying graph variables, solved as a fixed-lag sliding-window
smoother, matches or beats the online EKF+TL+NN on real data \emph{with no neural
network}.

We are explicit about what is and is not novel. For the linear--Gaussian INS
error model the batch maximum-a-posteriori (MAP) estimate on the factor-graph
chain is algebraically the Rauch--Tung--Striebel (RTS)
smoother~\cite{rauch1965}; ``FGO versus EKF'' as pure smoothing is therefore not
our contribution. Our contributions are:
\begin{enumerate}
\item An online aeromagnetic-compensation formulation in which the
time-varying TL coefficients are factor-graph variables of a fixed-lag
sliding-window smoother (iSAM2-style~\cite{kaess2012}), giving an NN-free
alternative to online EKF+TL+NN cold-start calibration
(Section~\ref{sec:method}).
\item A demonstration on real SGL~2020 data that this formulation matches or
beats a reproduced---not merely cited---online EKF+TL+NN on the primary
evaluation line, and generalizes across five lines, three flights, and two maps,
where the causal EKF instead diverges (Section~\ref{sec:results}).
\item Physical scalar-magnetometer error factors (heading error, dead-zone
weighting, hard-iron bias, drift) estimated jointly with navigation, and an
observability analysis of joint compensation and navigation that identifies an
expressiveness ``sweet spot'' and an exogeneity design rule
(Sections~\ref{sec:sensor} and~\ref{sec:discussion}).
\end{enumerate}

\section{Background and Problem Formulation}\label{sec:background}

\subsection{INS error state and map-matching measurement}
We use the standard 17-state Pinson INS error model augmented with a
first-order Gauss--Markov (FOGM) catch-all magnetic state, with error state
$\bx\in\mathbb{R}^{n}$ and discrete linearized dynamics
$\bx_{t+1}=\bF_t\,\bx_t+\mathbf{w}_t$, $\mathbf{w}_t\sim\mathcal{N}(0,\mathbf{Q}_t)$,
where $\bF_t$ is the Pinson transition matrix. The scalar magnetometer
measurement is modeled as
\begin{equation}
z_t = h_{\mathrm{map}}\!\left(\mathbf{p}_t\right) + c_t + \eta_t,
\qquad \eta_t\sim\mathcal{N}(0,R),
\label{eq:meas}
\end{equation}
where $h_{\mathrm{map}}$ interpolates the anomaly map (plus IGRF core) at position
$\mathbf{p}_t$, and $c_t$ is the aircraft interference to be compensated. A
position error $\delta\mathbf{p}$ perturbs the prediction by
$\nabla h_{\mathrm{map}}\cdot\delta\mathbf{p}$; this gradient is the sole source of
horizontal position information.

\subsection{Tolles--Lawson compensation}
The interference is modeled as $c_t=\bA_t^{\!\top}\boldsymbol{\beta}_t$, where
$\bA_t$ is the TL basis row (permanent, induced, and eddy-current terms, plus an
optional bias) formed from the fluxgate vector at time $t$~\cite{tolles1950}, and
$\boldsymbol{\beta}_t$ are the TL coefficients. Classically $\boldsymbol{\beta}$
is constant and fit offline; online/cold-start operation makes it a slowly
time-varying state.

\subsection{Batch MAP estimation as a factor graph}
Stacking the prior, the process factors, and the measurement factors over a window
of $N$ epochs, the MAP estimate minimizes
\begin{equation}
\begin{aligned}
J(\bx_{1:N})=\;&\lVert\bx_1-\bx_0\rVert^2_{\mathbf{P}_0^{-1}}
+\sum_{t=1}^{N-1}\lVert\bx_{t+1}-\bF_t\bx_t\rVert^2_{\mathbf{Q}_t^{-1}}\\
&+\sum_{t=1}^{N}\rho\!\left(\frac{z_t-h_t(\bx_t)}{\sqrt{R}}\right),
\end{aligned}
\label{eq:cost}
\end{equation}
with $\rho(\cdot)$ a (possibly robust) kernel. For $\rho(u)=u^2$ and the
linear--Gaussian chain the minimizer of~\eqref{eq:cost} equals the RTS
smoother~\cite{rauch1965}; we solve it either by an iterated RTS pass
(Gauss--Newton relinearization) or by a global sparse Gauss--Newton/QR
(square-root SAM~\cite{dellaert2017}) step, and use iteratively reweighted least
squares for robust $\rho$ (Huber/Cauchy).

\section{Method}\label{sec:method}

\subsection{Online TL coefficients as factor-graph variables}
We append the TL coefficients $\boldsymbol{\beta}$ to the error state, so the
measurement factor becomes
\begin{equation}
z_t = \bA_t^{\!\top}\boldsymbol{\beta}_t + h_{\mathrm{map}}(\mathbf{p}_t) +
S_t + \eta_t,
\end{equation}
with $S_t$ the FOGM catch-all. A prior factor anchors $\boldsymbol{\beta}$ and
process factors let it drift slowly. The measurement Jacobian row is
$\bH_t=[\,\partial h_{\mathrm{map}}/\partial\mathbf{p}\;\;\mathbf{0}\;\;\bA_t^{\!\top}\;\;1\,]$.
This is the batch analog of the online-TL EKF; because the batch solution uses the
\emph{entire} window, early navigation states benefit from calibration information
that only becomes observable after later maneuvers.

\subsection{Fixed-lag sliding-window smoother}
A single batch over a long line forces one static
$\boldsymbol{\beta}$, which underfits time-varying interference. We instead run
the estimator as an iSAM2-style fixed-lag smoother~\cite{kaess2012}: the flight is
processed in overlapping windows of length $L_w$ with overlap $L_o$; each window
is a batch FGO. A window commits only its leading stride $L_w-L_o$ epochs and
carries the full smoothed state and covariance forward as the prior for the next
window, so the calibration \emph{adapts} along the flight while navigation stays
continuous across windows. The trailing overlap acts as smoother look-ahead and is
re-committed by the next window. This is the factor-graph analog of the online NN
in~\cite{hager2026}: the per-window relinearization, not a learned residual,
supplies the needed time variation.

\subsection{Physical sensor-error factors}\label{sec:sensor}
Beyond TL, we cast physically-motivated scalar-magnetometer error terms as graph
variables estimated jointly with navigation: an optically-pumped-magnetometer
heading error as a Fourier series in the sensor--field angle $\theta$ (light shift
$\propto\cos\theta$, nonlinear Zeeman $\propto\cos 2\theta$), a fluxgate hard-iron
bias $\mathbf{m}\cdot\hat{\mathbf{u}}_{\mathrm{body}}$, a linear drift, and a
dead-zone heteroscedastic measurement weight $R/\lvert\sin 2\theta\rvert^2$. A
position-only grid or point-mass estimator structurally cannot carry such
sensor-error states; the graph formulation can.

\section{Reproduced Online EKF+TL+NN Baseline}\label{sec:baseline}
To compare against~\cite{hager2026} on equal footing we \emph{re-run}, rather than
merely cite, an online EKF with TL basis features and NN weights carried as EKF
states and learned online (the reference implementation of that filter family).
Reproducing the cold start required three fixes, each addressing a distinct
failure: (i) initializing the NN output bias to the onboard-estimated interference
DC (otherwise the first residual is the full interference offset and the gain
spikes); (ii) restricting the NN inputs to the attitude-derived TL basis and
\emph{excluding} the scalar magnetometer itself, which otherwise lets the network
subtract the very anomaly used for navigation; and (iii) de-trusting the map
through the cold-start transient. The resulting filter lands inside the paper's
published band (Section~\ref{sec:results}), and we use it as the comparison
baseline.

\section{Experiments}\label{sec:experiments}
All experiments use the public SGL~2020 dataset~\cite{gnadt2020} and the Ottawa
anomaly maps (Eastern\_395, Renfrew\_395), which download automatically; every
result is produced by continuous-integration scripts on the code branch. The
metric is horizontal distance-RMS (DRMS) after a \SI{10}{\minute} warm-up (the
convention of~\cite{hager2026}). Cold start means TL coefficients initialized to
zero with no calibration flight. Filters use identical process/measurement noise
across the EKF/FGO comparison so that only the estimator structure differs.

\section{Results}\label{sec:results}

\subsection{Batch FGO versus EKF on real data}
On line 1003.02 (Flt1003, Eastern\_395, \SI{63}{\minute}) with the compensated
stinger magnetometer, batch FGO improves horizontal DRMS from the EKF's
\SI{28.3}{\meter} to \SI{14.3}{\meter} (\SI{14.0}{\meter} with a Huber kernel),
against \SI{114.6}{\meter} free-inertial (Table~\ref{tab:flt1003}). On a
\SI{10}{\minute} segment the iterated-RTS and global sparse-GN/QR solvers agree
(\SI{18.4}{} vs \SI{19.4}{\meter}), confirming both reach the same MAP estimate.

\begin{table}[t]
\caption{Batch estimator on Flt1003 line 1003.02 (stinger Mag 1).}
\label{tab:flt1003}
\centering
\begin{tabular}{lr}
\toprule
Method & DRMS [m] \\
\midrule
INS (no aiding) & 114.6 \\
EKF & 28.3 \\
FGO (RTS) & 14.3 \\
FGO (RTS + Huber) & \textbf{14.0} \\
\bottomrule
\end{tabular}
\end{table}

\subsection{Cold-start line: FGO versus online EKF+TL+NN}
Table~\ref{tab:1007} reports the primary evaluation line 1007.06 (Flt1007,
Renfrew\_395, \SI{87}{\minute}, full length), on the two uncompensated cabin
magnetometers. A single static-TL batch underfits and degrades over the long line
(\SI{123.7}{}/\SI{68.1}{\meter}). The fixed-lag sliding window recovers it to
\SI{37.0}{}/\SI{15.1}{\meter} (\SI{5}{\minute} window), and with a Huber kernel to
\SI{32.6}{}/\SI{14.2}{\meter}. Our reproduced online EKF+TL+NN reaches
\SI{40.0}{}/\SI{17.5}{\meter}, inside the band published for that
filter~\cite{hager2026} (TL+NN \SI{37}{}/\SI{14}{\meter}; TL-only
\SI{58}{}/\SI{15}{\meter}). The sliding-window FGO thus matches or beats the
NN-augmented filter \emph{without any neural network}; the adaptive-TL
relinearization plays the role the online NN plays.

\begin{table}[t]
\caption{Cold start, line 1007.06, uncompensated cabin magnetometers.
DRMS [m] after \SI{10}{\minute} warm-up. INS \SI{318}{\meter};
EKF on compensated Mag 1 $\approx$ \SI{20}{\meter}.}
\label{tab:1007}
\centering
\begin{tabular}{lrr}
\toprule
Method (no NN unless noted) & Mag 4 & Mag 5 \\
\midrule
FGO-online, static batch TL & 123.7 & 68.1 \\
FGO-online, window \SI{5}{\minute} & 37.0 & 15.1 \\
FGO-online, window \SI{2}{\minute} & 45.9 & 17.1 \\
\textbf{FGO-online, window \SI{5}{\minute} + Huber} & \textbf{32.6} & \textbf{14.2} \\
\midrule
EKF+TL+NN, reproduced~\cite{hager2026} & 40.0 & 17.5 \\
EKF+TL+NN, published~\cite{hager2026} & 37 & 14 \\
EKF+TL-only, published~\cite{hager2026} & 58 & 15 \\
\bottomrule
\end{tabular}
\end{table}

\subsection{Breadth: the advantage generalizes}
To test that the result is not line-specific, Table~\ref{tab:breadth} runs the
\emph{same} cold-start TL model through the causal filter (online-TL EKF) and the
factor graph (\SI{5}{\minute} window + Huber) on five long survey lines spanning
three flights and two maps. The sliding-window FGO beats the causal EKF on
\num{8} of \num{9} cold-start cases (one EKF run failed off-map). The decisive
pattern is robustness: on the noisy Mag 4 the causal EKF \emph{diverges} from the
cold start to tens of kilometers on three lines
(\num{41626}, \num{35482}, \SI{17179}{\meter}), whereas the batch/window FGO stays
bounded (\SI{42.6}{}, \SI{38.6}{}, \SI{193.9}{\meter}). The fixed-lag smoother
re-linearizes over each window, so early navigation is protected by calibration
that only becomes observable later---something a one-pass causal filter cannot do.
On the cleaner Mag 5 the FGO is uniformly better, and it beats even the
compensated-stinger EKF on several lines. The lone exception (Flt1006 1006.08) is
a \SI{14}{\minute} line on which every method is weak, reflecting a
map-information floor for short lines rather than a filter failure.

\begin{table}[t]
\caption{Breadth: cold-start cabin magnetometers, DRMS [m] after
\SI{10}{\minute}. Same TL model; causal online-TL EKF vs \SI{5}{\minute}-window
FGO+Huber. ``div.''\ denotes divergence to $>\SI{10}{\kilo\meter}$.}
\label{tab:breadth}
\centering
\setlength{\tabcolsep}{4pt}
\begin{tabular}{llccrr}
\toprule
Flight & Line & Map & Mag & EKF-online & FGO-win \\
\midrule
Flt1003 & 1003.02 & E & 4 & div.\ (41626) & \textbf{42.6} \\
Flt1003 & 1003.02 & E & 5 & 28.1 & \textbf{21.7} \\
Flt1003 & 1003.08 & R & 4 & div.\ (fail) & \textbf{26.1} \\
Flt1003 & 1003.08 & R & 5 & 21.1 & \textbf{12.4} \\
Flt1006 & 1006.08 & E & 4 & div.\ (17179) & \textbf{193.9} \\
Flt1006 & 1006.08 & E & 5 & \textbf{117.5} & 122.0 \\
Flt1007 & 1007.02 & E & 4 & div.\ (35482) & \textbf{38.6} \\
Flt1007 & 1007.02 & E & 5 & 31.6 & \textbf{14.5} \\
Flt1007 & 1007.06 & R & 4 & 46.7 & \textbf{32.7} \\
Flt1007 & 1007.06 & R & 5 & 17.8 & \textbf{13.8} \\
\midrule
\multicolumn{4}{l}{FGO wins} & \multicolumn{2}{r}{\textbf{8 / 9}}\\
\bottomrule
\end{tabular}
\end{table}

\subsection{Sensor-error factors (simulation)}
On a simulated flight (Eastern\_395) with injected heading error, hard-iron bias,
drift, and dead-zone noise (INS reference \SI{31.1}{\meter}), cumulatively adding
the sensor-error factors reduces DRMS from \SI{36.6}{\meter} (no sensor states,
worse than INS because unmodeled error corrupts aiding) to \SI{4.8}{\meter} (full
model), a \SI{87}{\percent} reduction. Parameter recovery is clean for the drift
(\num{0.019} vs \num{0.020}\,\si{\nT\per\second}) and one bias component, but the
heading harmonics are observability-limited because the level-flight profile
sweeps the sensor--field angle by only $\approx 12^\circ$; this is an
observability property, not an estimator defect.

\section{Discussion}\label{sec:discussion}
\subsection{Why the factor graph helps}
The gain is not smoothing per se (which equals RTS in the linear--Gaussian case)
but the \emph{joint, windowed} estimation of a time-varying compensation with
navigation. A causal EKF must commit compensation and navigation online with no
future context; under a large uncompensated cold-start residual it attributes the
mismatch to position and diverges. The fixed-lag window defers commitment, letting
later-observable calibration correct earlier states---exactly the robustness seen
in Table~\ref{tab:breadth}.

\subsection{Observability of joint compensation and navigation}
Joint cold-start compensation and navigation is well-posed only within an
expressiveness \emph{sweet spot}: too weak a basis under-compensates and the
residual corrupts position (permanent-only TL diverged to
\SI{7631}{\meter} on 1007.06), while a basis that can reproduce the map signal
lets compensation absorb it and position drifts (observability collapse). We
explored scalar predictors of the collapse---the map-value confound $R^2$, its
cross-segment (out-of-sample) form, and a factor-graph map-\emph{gradient} confound
computed from the window Jacobian---and report honestly that none cleanly
separates safe from collapse-prone bases across estimators, and that a
gradient-gated smoother did not recover an injected collapse in the linear regime.
The robust, actionable finding is qualitative and is what every result here uses:
compensation should be driven by \emph{exogenous} (attitude/fluxgate,
position-independent) features, never by the scalar being navigated on. Turning
the sweet-spot boundaries into a quantitative identifiability (CRLB/PCRB)
condition is left as future work.

\subsection{Limitations}
Results use single-run DRMS on five lines; a journal-grade study adds Monte-Carlo
statistics and the full line set, a like-for-like grid/point-mass MMSE and MPF
baseline, and real (not simulated) sensor-error validation. Short lines remain
hard for all methods.

\section{Conclusion}
Casting online Tolles--Lawson aeromagnetic compensation as time-varying
factor-graph variables in a fixed-lag sliding-window smoother yields a
neural-network-free cold-start calibrator that, on real SGL~2020 data, matches or
beats a reproduced online EKF+TL+NN on the primary line and generalizes across
five lines, three flights, and two maps---most importantly by remaining bounded
where the causal EKF diverges. The factor graph's per-window relinearization
supplies the time-varying compensation that the competing method obtains from an
online neural network, at lower conceptual and computational complexity.

\begin{thebibliography}{1}
\bibitem{canciani2017} A.~Canciani and J.~Raquet, ``Airborne magnetic anomaly
navigation,'' \emph{IEEE Trans. Aerosp. Electron. Syst.}, vol.~53, no.~1,
pp.~67--80, 2017.
\bibitem{gnadt2020} A.~R. Gnadt \emph{et al.}, ``Signal enhancement for magnetic
navigation challenge problem,'' arXiv:2007.12158, 2020; SGL 2020 dataset,
MagNav.jl.
\bibitem{tolles1950} W.~E. Tolles and J.~D. Lawson, ``Magnetic compensation of
MH-1A aircraft,'' Airborne Instruments Lab., Rep.~1, 1950.
\bibitem{hager2026} Hager \emph{et al.}, ``Airborne MagNav with NN-augmented
online calibration,'' arXiv:2603.08265, 2026.
\bibitem{rauch1965} H.~E. Rauch, F.~Tung, and C.~T. Striebel, ``Maximum likelihood
estimates of linear dynamic systems,'' \emph{AIAA J.}, vol.~3, no.~8,
pp.~1445--1450, 1965.
\bibitem{dellaert2017} F.~Dellaert and M.~Kaess, ``Factor graphs for robot
perception,'' \emph{Found. Trends Robot.}, vol.~6, no.~1--2, pp.~1--139, 2017.
\bibitem{kaess2012} M.~Kaess \emph{et al.}, ``iSAM2: Incremental smoothing and
mapping using the Bayes tree,'' \emph{Int. J. Robot. Res.}, vol.~31, no.~2,
pp.~216--235, 2012.
\end{thebibliography}

\end{document}
